% !TEX program = xelatex
% Example: Figure float with caption and label
\documentclass[11pt, a4paper]{article}
\usepackage{graphicx}
\usepackage{float}

\begin{document}

\section*{العناصر الطافية}

% عنصر طافٍ في أعلى الصفحة
\begin{figure}[htbp]
  \centering
  \includegraphics[width=0.7\textwidth]{example-image}
  \caption{رسمٌ توضيحيٌّ يُمثّلُ العلاقةَ بينَ المتغيّرَين}
  \label{fig:example}
\end{figure}

% إحالة للشكل في النص
كما يَظهرُ في الشكلِ \ref{fig:example}، فإنَّ العلاقةَ خطّيةٌ.

% عنصر طافٍ في موضعٍ محدّدٍ بالقوة
\begin{figure}[H]
  \centering
  \includegraphics[width=0.5\textwidth]{example-image}
  \caption{هذا الشكلُ يَظهرُ في موضعِهِ المحدّد}
  \label{fig:forced}
\end{figure}

\end{document}
